% =====================================================================
%  Capítulo 8 — El agresor sexual: prevención y respuesta
% =====================================================================
\chapter{El agresor sexual: prevención y respuesta}

Para enfrentarse a un agresor sexual es aplicable todo lo visto en los puntos precedentes, pero hay que mencionar variables propias de este tipo de agresiones. La primera y más importante, confirmada por todas las estadísticas (incluida la Macroencuesta de Violencia contra la Mujer del Ministerio de Igualdad): en la gran mayoría de los casos el agresor NO es un extraño, sino alguien del entorno: pareja o expareja, conocidos, compañeros de trabajo o estudios. El estereotipo del desconocido en el callejón existe, pero es minoritario.

\section{Prevención}

\begin{itemize}
  \item Es posible cortar muchas situaciones poniendo límites enérgicos en la etapa previa, cuando se hace evidente la presión psicológica: nombrar la conducta (\enquote{esto es acoso}), exigir que cese y, si procede, amenazar con denuncia delante de testigos. El acosador cuenta con el silencio.
  \item Evita compartir lugares de riesgo con un potencial agresor: ascensores, \enquote{entrevistas de trabajo} fuera de la oficina y del horario laboral, quedadas a solas con desconocidos sin que nadie sepa dónde estás. Este tipo de ataques solo se realizan en lugares seguros para el atacante: niégale ese terreno.
  \item No des información personal (teléfono, dirección, rutinas) a cualquier compañero o jefe solo por agradar.
  \item En citas con personas conocidas por internet: primer encuentro en lugar público, avisa a alguien de confianza (o usa la función Guardián/Acompáñame de AlertCops), vigila tu bebida y ten transporte de vuelta propio.
  \item La sumisión química existe: no aceptes bebidas que no hayas visto servir y no dejes tu vaso sin vigilancia. Si notas mareo o confusión desproporcionados, pide ayuda inmediatamente a personal del local o amigas.
\end{itemize}

\section{Si el ataque llega: qué dice la evidencia}

Durante décadas se aconsejó de forma genérica \enquote{no resistirse}. La investigación actual lo matiza de forma importante: los estudios sobre resistencia \citep{ullman1997, ullman2007, tark2014} muestran que la resistencia activa ---gritar, forcejear, golpear, huir--- reduce significativamente la probabilidad de que la violación se consume, y que en la mayoría de los casos no incrementa las lesiones físicas adicionales: la lesión suele preceder a la resistencia, no ser su consecuencia. Por eso los programas modernos de autodefensa femenina entrenan la resistencia verbal y física inmediata y decidida.

Dicho esto, ninguna estadística decide por ti en el momento: cada situación es única (armas, encierro, varios agresores) y TODA estrategia de supervivencia es legítima, incluida la sumisión táctica para ganar tiempo y esperar una oportunidad. Quien sobrevive, lo hizo bien. La culpa es siempre del agresor.

\section{La oportunidad y la sorpresa}

Comúnmente los agresores sexuales no tienen especial experiencia en combate, por lo que una persona decidida puede causar un gran daño en un individuo ocupado en otros menesteres que no son la pelea. El agresor suele subestimar a la mujer porque la supone débil: si aparentas rendición, bajará su atención defensiva, y ese es el momento de actuar con total decisión. El riesgo psicológico a evitar es ponerse mentalmente en un plano de inferioridad que anule la acción: contra eso se entrena.

Una vez lograda la oportunidad, se imponen procedimientos que deben repasarse mentalmente y que, aunque parezcan excesivos, son proporcionados frente a una violación:

\begin{itemize}
  \item Con la cara del agresor cerca, hundir fuerte y profundamente los pulgares en los ojos.
  \item Empuñar elementos del entorno: en una oficina, un bolígrafo (asegurando la parte trasera con el pulgar) contra las partes blandas.
  \item La mandíbula desarrolla una potencia enorme: usar la boca para morder y desgarrar partes blandas de cara y cuello.
  \item Combinar todo con golpes contundentes de codo o estrellando la cabeza del agresor contra elementos duros, gritar para alertar, y retirarse rápidamente para pedir ayuda y denunciar.
\end{itemize}

Nunca sabemos las intenciones finales del agresor, así que no hay que tener ningún reparo ni compasión a la hora de defenderse.

\section{Si un asaltante te agarra, recuerda}

\begin{itemize}
  \item Si te agarra bien, será mucho más difícil escapar: actúa pronto.
  \item Si sus manos están ocupadas agarrándote, otros objetivos (ojos, garganta, genitales) se presentan por sí mismos.
  \item Muévete: hay que ser escurridiza.
  \item Grita órdenes concretas: \enquote{¡Suéltame!}, \enquote{¡Fuego!}, \enquote{¡Llamad a la policía!}. El grito te activa, te hace respirar, desconcierta y atrae testigos.
\end{itemize}

Si intentas una técnica para escapar y no funciona, no te rindas: intenta otra. Haz lo inesperado. Trabajar con un instructor que lleva protecciones te dará una idea de cómo se siente realmente intentar liberarse.
